% GENERATED FILE. Edit canonical lessons, then run npm run generate:books.
% canonical-source-hash: fnv1a64:821183f20abf042c
% canonical-lessons: RU-C504-shashki, RU-C504-peshka, RU-C504-kubik, RU-C504-koloda, RU-C504-zhulnichat

\chapter{Draughts, Dice and Cards}
\label{ch:ru-draughts-dice-and-cards}
\hlchaptermodality{504}

\begin{chapteropening}
\textbf{By the end of this chapter:} \emph{I can talk about board games and card games.}

Complete the last lesson of chapter 504: I can talk about board games and card games.
\end{chapteropening}

\section[\texorpdfstring{\ru{шашки} (sháshki)}{шашки (sháshki)}]{\ru{шашки} (sháshki) — draughts, checkers}

\label{lesson:RU-C504-shashki}



\begin{quote}
Before the new one: say the Russian for a crossword, then the Russian for to solve.
\end{quote}

\subsection*{You'll want to know: \ru{шашки}}
\textbf{\ru{шашки}} (\emph{sháshki}) — \textquotedblleft{}draughts, checkers\textquotedblright{}.

\begin{quote}
\ru{По} \ru{вечерам} \ru{дед} \ru{играет} \ru{с} \ru{соседом} \ru{в} \ru{шашки}.
\end{quote}

\begin{quote}
\emph{Po vecherám ded igráyet s sosédom v sháshki.}
\end{quote}

\begin{quote}
In the evenings Grandad plays draughts with the neighbour.
\end{quote}

Plural. \emph{\ru{Играть} \ru{в} \ru{шашки}} — to play draughts.

\begin{cousinweb}
EtymDB lists \emph{\ru{шашки}} without a source; it is usually traced to \emph{\ru{шах}} 'check', from Persian \emph{shāh} 'king', like \emph{\ru{шахматы}} 'chess'.
\end{cousinweb}

\begin{grammarlens}[title={What you've built}]
One more word for telling, explaining and describing at B1.
\end{grammarlens}

\subsection*{Your turn}
\begin{itemize}
\raggedright
  \item \emph{Say it:} \emph{sháshki}
  \item \emph{Say it:} \emph{sháshki}, once more
  \item \emph{Say it:} a sentence of your own with \emph{sháshki}
  \item \emph{Recall:} say the Russian for a crossword, then the Russian for to solve, then say \emph{sháshki} again
\end{itemize}

\begin{tcolorbox}[breakable,colback=teal!4,colframe=teal!35!black,title={Before you move on}]
What does \ru{шашки} mean? (\textquotedblleft{}Draughts, checkers\textquotedblright{}.) Say it once more.
\end{tcolorbox}
\section[\texorpdfstring{\ru{пешка} (péshka)}{пешка (péshka)}]{\ru{пешка} (péshka) — a pawn}

\label{lesson:RU-C504-peshka}



\begin{quote}
Before the new one: say the Russian for to solve, then the Russian for draughts.
\end{quote}

\subsection*{You'll want to know: \ru{пешка}}
\textbf{\ru{пешка}} (\emph{péshka}) — \textquotedblleft{}a pawn\textquotedblright{}.

\begin{quote}
\ru{Пешка} — \ru{самая} \ru{слабая} \ru{фигура} \ru{в} \ru{шахматах}.
\end{quote}

\begin{quote}
\emph{Péshka — sámaya slábaya figúra v shákhmatakh.}
\end{quote}

\begin{quote}
The pawn is the weakest piece in chess.
\end{quote}

Feminine. Figuratively, a person of no importance.

\begin{cousinweb}
\emph{\ru{пешка}} is built on \emph{\ru{пеший}} 'on foot', so a pawn is a foot soldier; EtymDB lists \emph{\ru{пешка}} without a source.
\end{cousinweb}

\begin{grammarlens}[title={What you've built}]
One more word for telling, explaining and describing at B1.
\end{grammarlens}

\subsection*{Your turn}
\begin{itemize}
\raggedright
  \item \emph{Say it:} \emph{péshka}
  \item \emph{Say it:} \emph{péshka}, once more
  \item \emph{Say it:} a sentence of your own with \emph{péshka}
  \item \emph{Recall:} say the Russian for to solve, then the Russian for draughts, then say \emph{péshka} again
\end{itemize}

\begin{tcolorbox}[breakable,colback=teal!4,colframe=teal!35!black,title={Before you move on}]
What does \ru{пешка} mean? (\textquotedblleft{}A pawn\textquotedblright{}.) Say it once more.
\end{tcolorbox}
\section[\texorpdfstring{\ru{кубик} (kúbik)}{кубик (kúbik)}]{\ru{кубик} (kúbik) — a die; a small cube}

\label{lesson:RU-C504-kubik}



\begin{quote}
Before the new one: say the Russian for draughts, then the Russian for a pawn.
\end{quote}

\subsection*{You'll want to know: \ru{кубик}}
\textbf{\ru{кубик}} (\emph{kúbik}) — \textquotedblleft{}a die; a small cube\textquotedblright{}.

\begin{quote}
\ru{Он} \ru{бросил} \ru{кубик}, \ru{и} \ru{выпала} \ru{шестёрка}.
\end{quote}

\begin{quote}
\emph{On brósil kúbik, i výpala shestyórka.}
\end{quote}

\begin{quote}
He threw the die and got a six.
\end{quote}

Masculine. A diminutive of \emph{\ru{куб}} — a cube. \emph{\ru{Кубики}} are also a child's building blocks.

\begin{cousinweb}
\emph{\ru{кубик}} is a diminutive of \emph{\ru{куб}} 'cube', which EtymDB traces through German \emph{Kubus} and Latin \emph{cubus} to Greek \emph{kýbos} 'a cube'; EtymDB has no entry for \emph{\ru{кубик}}.
\end{cousinweb}

\begin{grammarlens}[title={What you've built}]
One more word for telling, explaining and describing at B1.
\end{grammarlens}

\subsection*{Your turn}
\begin{itemize}
\raggedright
  \item \emph{Say it:} \emph{kúbik}
  \item \emph{Say it:} \emph{kúbik}, once more
  \item \emph{Say it:} a sentence of your own with \emph{kúbik}
  \item \emph{Recall:} say the Russian for draughts, then the Russian for a pawn, then say \emph{kúbik} again
\end{itemize}

\begin{tcolorbox}[breakable,colback=teal!4,colframe=teal!35!black,title={Before you move on}]
What does \ru{кубик} mean? (\textquotedblleft{}A die; a small cube\textquotedblright{}.) Say it once more.
\end{tcolorbox}
\section[\texorpdfstring{\ru{колода} (kolóda)}{колода (kolóda)}]{\ru{колода} (kolóda) — a pack of cards; a log}

\label{lesson:RU-C504-koloda}



\begin{quote}
Before the new one: say the Russian for a pawn, then the Russian for a die.
\end{quote}

\subsection*{You'll want to know: \ru{колода}}
\textbf{\ru{колода}} (\emph{kolóda}) — \textquotedblleft{}a pack of cards; a log\textquotedblright{}.

\begin{quote}
\ru{В} \ru{колоде} \ru{не} \ru{хватает} \ru{одной} \ru{карты}.
\end{quote}

\begin{quote}
\emph{V kolóde ne khvatáyet odnóy kárty.}
\end{quote}

\begin{quote}
One card is missing from the pack.
\end{quote}

Feminine. Its older sense is a log or a block of wood.

\begin{cousinweb}
EtymDB lists \emph{\ru{колода}} without a source; it is usually traced to Proto-Slavic \emph{kolda} 'a log, a block'.
\end{cousinweb}

\begin{grammarlens}[title={What you've built}]
One more word for telling, explaining and describing at B1.
\end{grammarlens}

\subsection*{Your turn}
\begin{itemize}
\raggedright
  \item \emph{Say it:} \emph{kolóda}
  \item \emph{Say it:} \emph{kolóda}, once more
  \item \emph{Say it:} a sentence of your own with \emph{kolóda}
  \item \emph{Recall:} say the Russian for a pawn, then the Russian for a die, then say \emph{kolóda} again
\end{itemize}

\begin{tcolorbox}[breakable,colback=teal!4,colframe=teal!35!black,title={Before you move on}]
What does \ru{колода} mean? (\textquotedblleft{}A pack of cards; a log\textquotedblright{}.) Say it once more.
\end{tcolorbox}
\section[\texorpdfstring{\ru{жульничать} (zhúlnichat)}{жульничать (zhúlnichat)}]{\ru{жульничать} (zhúlnichat) — to cheat (at a game)}

\label{lesson:RU-C504-zhulnichat}



\begin{quote}
Before the new one: say the Russian for a die, then the Russian for a pack of cards.
\end{quote}

\subsection*{You'll want to know: \ru{жульничать}}
\textbf{\ru{жульничать}} (\emph{zhúlnichat}) — \textquotedblleft{}to cheat (at a game)\textquotedblright{}.

\begin{quote}
\ru{С} \ru{ним} \ru{нельзя} \ru{играть} \ru{в} \ru{карты}: \ru{он} \ru{жульничает}.
\end{quote}

\begin{quote}
\emph{S nim nelzyá igrát v kárty: on zhúlnichayet.}
\end{quote}

\begin{quote}
You can't play cards with him: he cheats.
\end{quote}

Imperfective, colloquial; present \ru{жульничаю}. Built on \emph{\ru{жулик}} — a crook.

\begin{cousinweb}
\emph{\ru{жульничать}} is built on \emph{\ru{жулик}} 'a crook'; EtymDB lists the verb without a source, and the origin of \emph{\ru{жулик}} is not settled.
\end{cousinweb}

\begin{grammarlens}[title={What you've built}]
One more word for telling, explaining and describing at B1.
\end{grammarlens}

\subsection*{Your turn}
\begin{itemize}
\raggedright
  \item \emph{Say it:} \emph{zhúlnichat}
  \item \emph{Say it:} \emph{zhúlnichat}, once more
  \item \emph{Say it:} a sentence of your own with \emph{zhúlnichat}
  \item \emph{Recall:} say the Russian for a die, then the Russian for a pack of cards, then say \emph{zhúlnichat} again
\end{itemize}

\begin{tcolorbox}[breakable,colback=teal!4,colframe=teal!35!black,title={Before you move on}]
What does \ru{жульничать} mean? (\textquotedblleft{}To cheat (at a game)\textquotedblright{}.) Say it once more.
\end{tcolorbox}
